% Add this to your document preamble if not already present:
% \usepackage{tabularx}
% \usepackage[table]{xcolor}

\chapter{Appendix}

\section*{Important Formulas and Definitions}

\begin{itemize}
  \item \textbf{Place Value:} Value of a digit = digit $\times$ place value.
  \item \textbf{Fraction to Decimal:} $\frac{a}{b} = a \div b$.
  \item \textbf{Prime Number:} A number greater than $1$ with exactly two factors: $1$ and itself.
  \item \textbf{Composite Number:} A number with more than two factors.
  \item \textbf{GCF:} Product of lowest powers of common primes.
  \item \textbf{LCM:} Product of highest powers of all primes present.
\end{itemize}

\section*{Further Reading and Resources}

\begin{itemize}
  \item "Elementary Mathematics" by Liping Ma
  \item "How to Be a Math Genius" by Mike Goldsmith
  \item Khan Academy: https://www.khanacademy.org/
  \item Purplemath: https://www.purplemath.com/
\end{itemize}

\newpage

\section*{Summary of Key Concepts}
\begin{table}[h!]
\centering
\renewcommand{\arraystretch}{1.3}
\begin{tabularx}{\textwidth}{|l|X|}
\hline
\rowcolor[HTML]{E0E0E0}
\textbf{Topic} & \textbf{Key Points} \\
\hline
Place Value & Each digit's value depends on its position. Example: $4275 = 4000 + 200 + 70 + 5$. \\
\hline
Decimals and Fractions & Decimals and fractions both represent parts of a whole. Example: $0.75 = \frac{3}{4}$. \\
\hline
Comparing and Ordering Numbers & Use $>$, $<$, $=$ to compare. Convert to decimals for easy comparison. Example: $\frac{2}{5} < \frac{3}{7}$. \\
\hline
Factors and Multiples & Factors divide a number evenly; multiples are products of a number and an integer. Example: Factors of $12$ are $1,2,3,4,6,12$. \\
\hline
Prime and Composite Numbers & Primes have exactly two factors (1 and itself); composites have more. Example: $2,3,5$ are prime; $4,6,8$ are composite. \\
\hline
GCF and LCM & GCF: largest common factor; LCM: smallest common multiple. Use prime factorization. Example: GCF of $12$ and $18$ is $6$; LCM is $36$. \\
\hline
\end{tabularx}
\caption{\textbf{Summary of Key Concepts}}
\end{table}


% Add any references or credits here.